\documentclass[11pt,a4paper]{article}

\usepackage[T1]{fontenc}
\usepackage[utf8]{inputenc}
\usepackage[margin=2.3cm]{geometry}
\usepackage{graphicx}
\usepackage{float}
\usepackage{xcolor}
\usepackage{amsmath}
\usepackage{caption}
\usepackage{parskip}
\usepackage[hidelinks]{hyperref}

\captionsetup{font=small,labelfont=bf,width=0.95\textwidth}
\setlength{\parskip}{0.6em}

\definecolor{boxgrey}{HTML}{F2F5F8}
\definecolor{inkblue}{HTML}{12263A}

\title{\vspace{-1.2cm}\textbf{Power side-channel input recovery on the CW305 board\\
\large Summary of the catch-up report}}
\author{Tarek Alhafez}
\date{9 September 2026\\
\small(condensed from the full catch-up report and its follow-up work)}

\begin{document}
\maketitle
\vspace{-0.8cm}

\section*{What this document is}

The short version of the catch-up report. It keeps the goal, the method, the figures and
the results, and drops the long reasoning. Where the full report later corrected a
number, only the corrected number appears here. References like ``(full report \S8.3)''
point into the long version.

\section*{1. Goal of the experiment}

I am reproducing the attack of Wei et al., \emph{I Know What You See} (ACSAC 2018). The
first convolution layer of a binarised neural network accelerator processes the input
image window by window. The power drawn by the chip depends on the pixels inside the
convolution window at that moment, so an attacker who measures only the current can
recover the input image. The attacker never sees the image.

The paper used a SAKURA-G board with a Spartan-6 FPGA and a 2.5\,GS/s oscilloscope. I do
not have that equipment. I use a CW305 board with an Artix-7 XC7A100T and a
ChipWhisperer-Lite. Two questions drive the project: whether the attack still works on
this cheaper chain, and what happens when the profiling conditions and the attack
conditions differ. The paper never reports the second one.

\begin{figure}[H]
  \centering
  \includegraphics[width=\linewidth]{figures/fig01_signal_chain.png}
  \caption{The measurement chain. The convolution runs on the FPGA. The PC only writes
  the input, starts the operation, and does the analysis afterwards. The power signal
  never passes through the PC before the ADC.}
  \label{fig:chain}
\end{figure}

The FPGA computes the convolution, not the host. Before every capture the script writes
one image and one kernel to the board, reads back all 676 output values and compares
them with a CPU model. If one value differs, the script refuses the capture. Every
capture used here came back $676/676$ bit-exact.

\section*{2. What was blocking the project}

\begin{sloppypar}
For weeks the detection looked weak and unstable. Two older files in the repository,
\texttt{FINAL\_RESULTS.md} and \texttt{REPORT\_1}, blame per-cycle extraction and
smearing. That explanation is wrong, and anyone who reads it first will optimise the
wrong thing.
\end{sloppypar}

The real cause was simpler. The board was still answering as the \emph{stock AES
example}: CW305 register pages 2, 3 and 4 returned \texttt{0x02}, \texttt{0x05},
\texttt{0x2e}, the AES signature. My bitstream had built correctly and met timing, but it
was not the image running on the FPGA. So the host wrote MNIST pixels and convolution
kernels into AES registers, and every captured trace held the power of an AES
encryption. The reconstruction was searching for image structure in traces that contained
no image.

I added an authenticity gate that reads those pages before any capture and refuses to run
if the AES signature is there. The gate earns its place because it can fail: it blocked
every capture attempt until the correct bitstream was on the board, and passed afterwards.

\begin{figure}[H]
  \centering
  \includegraphics[width=\linewidth]{figures/fig02_authenticity_gate_nodates.png}
  \caption{The same gate, before and after the bitstream was fixed. A test that never
  fails proves nothing. This one refused the AES session and allowed the real one.}
  \label{fig:gate}
\end{figure}

\section*{3. Method}

\subsection*{3.1 The signal}

The FPGA runs at 5\,MHz and the ChipWhisperer ADC takes four samples per FPGA clock,
locked to that clock. Each convolution window takes 8 clock cycles, so 32 ADC samples.
The 676 windows give 21632 useful samples out of the 22144 captured.

\begin{figure}[H]
  \centering
  \includegraphics[width=\linewidth]{figures/fig03_raw_trace.png}
  \caption{Real captured power signal for one probe kernel and one image. Top: all 676
  windows, drawn as the envelope of the 32 samples inside each window. Bottom left: the
  raw samples of four windows, where the window boundaries are visible. Bottom right:
  the single number that the attack keeps for each window.}
  \label{fig:raw}
\end{figure}

My chain has one advantage over the paper's. Their oscilloscope sampled asynchronously,
which smeared the per-cycle power and cost them a whole section (their Section 5) on DC
restoration and per-cycle curve fitting. My ADC is phase-locked to the FPGA clock, so I
read per-cycle power directly. I did not reproduce their Section 5 and I do not need it.
Skipping it discards no information the attack uses.

From each trace I keep one feature per window: the mean absolute value of its 32 samples.
Nine probe kernels give a 9-number vector per window. The digit shows up in that feature
before any attack algorithm runs, which is where I first believed the leakage was real.

\begin{figure}[H]
  \centering
  \includegraphics[width=\linewidth]{figures/fig04_feature_map.png}
  \caption{From the captured signal to the attack feature. The FPGA receives a
  binarised image, not the grey image. The power feature map of a single kernel already
  shows the shape of the digit.}
  \label{fig:feat}
\end{figure}

\subsection*{3.2 The attack}

The attack has two stages.

\emph{Profiling.} I use images I know. For every convolution window I store the 9-kernel
power vector together with the true $3\times3$ binary patch that produced it. There are
$2^9 = 512$ possible patches. Forty profiling images give a template of
$40 \times 676 = 27040$ rows.

\emph{Attack.} Here I have only traces. The ground truth never enters the attack bundle:
the code refuses to run if the bundle contains an \texttt{image} or \texttt{label} array.
For each window I split the 9 kernels into 3 groups of 3, search each group in the
template with a kD-tree inside radius $\delta$, and intersect the three result sets to get
candidate patches. A greedy selection, following Algorithm 2 of the paper, picks one
candidate per window by preferring the one that agrees with its neighbours. The
overlapping patches then vote on each pixel.

\begin{figure}[H]
  \centering
  \includegraphics[width=\linewidth]{figures/fig05_algorithm.png}
  \caption{The active power-template attack as implemented here. The grouping, $\delta$
  and the number of probe kernels are the same as in the paper.}
  \label{fig:algo}
\end{figure}

\subsection*{3.3 Where this differs from the paper}

I did not reproduce the paper exactly. Three differences change how the results should be
read.

\begin{itemize}
\item \textbf{Binary pixels.} In the paper the first layer takes real 0--255 pixels
      against binary weights, so the power depends on the grey level and the attacker can
      recover a grey image. My binary design thresholds the pixel at 127 before the
      convolution, which caps the best possible result at a one-bit silhouette. I
      therefore score against the \emph{binarised} image and report F1, IoU and MSSIM
      instead of pixel accuracy alone. A separate grey-pixel design lifts this limit and
      supplies the generative results in Section 4.
\item \textbf{The baseline is always printed.} Most MNIST pixels are background, so pixel
      accuracy looks good even when the attack recovers nothing. Every table prints the
      accuracy of the trivial ``everything is background'' answer beside it.
\item \textbf{Geometry.} The paper uses $28\times28$ ``same'' padding, 784 cycles. My
      hardware does a valid convolution: $26\times26$, 676 windows. The window-to-pixel
      mapping that the attack depends on works the same way in both.
\end{itemize}

\begin{figure}[H]
  \centering
  \includegraphics[width=\linewidth]{figures/fig06_geometry.png}
  \caption{The geometry of my layer is not the same as the paper's layer, and my layer
  binarises the pixel before the convolution.}
  \label{fig:geom}
\end{figure}

\begin{figure}[H]
  \centering
  \includegraphics[width=\linewidth]{figures/fig12_comparison.png}
  \caption{My setup next to the setup of the paper.}
  \label{fig:cmp}
\end{figure}

\section*{4. Results}

\begin{table}[H]
  \centering
  \small
  \resizebox{\textwidth}{!}{\begin{tabular}{lrrrrrrrr}
\hline
Run & Pixel acc. & All-background & F1 & IoU & MSSIM & Pixel dist. & Recog. & $n$ \\
\hline
300 profile images, control evaluation set & 0.953 & 0.887 & 0.759 & 0.612 & 0.713 & 11.9 & 0.93 & 40 \\
40 profile / 40 evaluation, fresh images & 0.968 & 0.887 & 0.843 & 0.728 & 0.816 & 8.2 & 1.00 & 40 \\
40 profile / 40 evaluation, earlier run & 0.960 & 0.888 & 0.792 & 0.656 & 0.769 & 10.2 & 0.90 & 40 \\
Probe mismatch (one-hot template vs trained traces) & 0.811 & 0.888 & 0.119 & 0.063 & 0.094 & 48.1 & 0.20 & 40 \\
Single hardware repeat (no averaging) & 0.919 & 0.888 & 0.494 & 0.328 & 0.501 & 20.6 & 0.65 & 40 \\
\hline
\end{tabular}
}
  \caption{Main results. Every row is a real hardware capture. Compare the
  ``all-background'' column with the pixel accuracy on its left.}
  \label{tab:headline}
\end{table}

\begin{table}[H]
  \centering
  \small
  \resizebox{\textwidth}{!}{\begin{tabular}{llrrccc}
\hline
Capture & Role & Clock & Gain & Not stock AES & Functional check & Mism. \\
\hline
rerun\_20260830\_trained\_80 & profiling (all condition tests) & 5.0\,MHz & 40\,dB & yes & 676/676 & 0 \\
rerun\_20260830\_onehot\_80 & profiling (probe mismatch) & 5.0\,MHz & 40\,dB & yes & 676/676 & 0 \\
prof\_5m & profiling (clock mixing, 5 MHz) & 5.0\,MHz & 40\,dB & yes & 676/676 & 0 \\
prof\_10m & profiling (clock mixing, 10 MHz) & 10.0\,MHz & 40\,dB & yes & 676/676 & 0 \\
grey\_p2p\_2000 & grey design, generative reconstruction & 5.0\,MHz & 40\,dB & yes & 676/676 & 0 \\
hard\_C1\_control & attack & 5.0\,MHz & 40\,dB & yes & 676/676 & 0 \\
hard\_C2\_reprogram & attack & 5.0\,MHz & 40\,dB & yes & 676/676 & 0 \\
hard\_C3\_clk7m5 & attack & 7.5\,MHz & 40\,dB & yes & 676/676 & 0 \\
hard\_C4\_clk10m & attack & 10.0\,MHz & 40\,dB & yes & 676/676 & 0 \\
hard\_C5\_gain30 & attack & 5.0\,MHz & 30\,dB & yes & 676/676 & 0 \\
hard\_C6\_gain50 & attack & 5.0\,MHz & 50\,dB & yes & 676/676 & 0 \\
paperscale\_500 & profiling only (300-image template) & 5.0\,MHz & 40\,dB & yes & 676/676 & 0 \\
\hline
\end{tabular}
}
  \caption{Hardware provenance of the captures used here. The functional check writes an
  image and a kernel to the FPGA, reads back the 676 outputs and compares them with a CPU
  model.}
\end{table}

\subsection*{4.1 The attack works on real hardware}

With the search radius tuned, foreground F1 reaches $\mathbf{0.969}$ and the digit
recogniser reads all 40 test images correctly. No part of the measurement is simulated.

The $0.843$ that runs through the older text came from the radius I carried over from the
paper, which sits far off the optimum for this hardware. Radius and template size have to
be tuned together. At a fixed radius, going from 40 to 300 profiling images made the
attack \emph{worse}; tuned, the larger template wins, $0.969$ against $0.958$. The best
radius shrinks as the template grows (full report \S8.1).

\begin{table}[H]
  \centering
  \small
  \resizebox{\textwidth}{!}{\begin{tabular}{lrrrrrrrrr}
\hline
Template & $\delta$ = 1.4 & $\delta$ = 1 & $\delta$ = 0.6 & $\delta$ = 0.4 & $\delta$ = 0.25 & $\delta$ = 0.15 & $\delta$ = 0.08 & $\delta$ = 0.06 & $\delta$ = 0.03 \\
\hline
40 profiling images & 0.730 & 0.843 & 0.912 & 0.940 & \textbf{0.958} & 0.936 & 0.865 & 0.820 & 0.736 \\
300 profiling images & 0.696 & 0.759 & 0.849 & 0.890 & 0.926 & 0.949 & 0.967 & \textbf{0.969} & 0.946 \\
\hline
\end{tabular}
}
  \caption{Foreground F1 against the search radius $\delta$. Both curves are unimodal,
  and the value carried over from the paper sits far off the optimum for this hardware.}
  \label{tab:delta}
\end{table}

\subsection*{4.2 Averaging is the main practical cost}

The convolution unit is small next to the rest of the fabric, so everything else on the
chip dilutes its power. That dilution is the limit on this board. With one hardware
repeat the attack sits barely above the trivial baseline. F1 and MSSIM move with
averaging; pixel accuracy hardly does.

The sweep is a fair one. A single capture stored the individual repeats, and each point
uses only the first $R$ of them, so images, temperature and gain stay identical across the
sweep and the averaging is the only thing that changes.

\begin{figure}[H]
  \centering
  \includegraphics[width=\linewidth]{figures/fig07_repeat_sweep.png}
  \caption{The averaging sweep. The left plot shows why the baseline must be drawn:
  pixel accuracy is already high at one repeat, but it sits close to the trivial answer.
  The metrics that move are F1 and MSSIM.}
  \label{fig:sweep}
\end{figure}

\subsection*{4.3 What survives a change of conditions, and what does not}

I built one template from 40 images at 5\,MHz and 40\,dB gain, then attacked the
\emph{same 40 new images} recaptured under several conditions. I checked the evaluation
sets with SHA-256 over the stacked ground truth and they all hash the same, so differences
between rows do not come from the images. Every condition passed the functional check
bit-exact, including 10\,MHz.

\begin{figure}[H]
  \centering
  \includegraphics[width=\linewidth]{figures/fig08_conditions.png}
  \caption{Template transfer under changed conditions. One profile, six capture
  conditions, the same 40 images.}
  \label{fig:cond}
\end{figure}

\begin{figure}[H]
  \centering
  \includegraphics[width=0.72\linewidth]{figures/fig09_examples.png}
  \caption{The same six images, recovered under different conditions. The empty row for
  the 10\,MHz capture is not a drawing mistake: at that clock the reconstruction returns
  an all-background image for every one of the 40 test images.}
  \label{fig:ex}
\end{figure}

\begin{itemize}
\item \textbf{Amplifier gain does not matter.} I changed it in both directions, to
      30\,dB and 50\,dB, and F1 held at $0.832$ and $0.831$ against $0.843$ for the
      control. The feature extraction normalises each trace by its own median and MAD, so
      a gain change amounts to a change of scale and the normalisation removes it.
\item \textbf{Clock does matter.} A clock change alters the shape of the power inside the
      window, and no normalisation undoes that. At 7.5\,MHz F1 falls to about $0.27$. At
      10\,MHz, attacked with a 5\,MHz template, the reconstruction returns all-background
      for every test image, and pixel accuracy lands on $0.887$, exactly the trivial
      baseline. This is the clearest case in the report for printing that column.
\item \textbf{The ``FPGA reprogrammed'' row is withdrawn.} It reported F1 $0.818$ and I
      first read it as the template surviving reconfiguration. Reprogramming a CW305 that
      is already configured does nothing, so that capture almost certainly ran on the same
      configuration as the control. The number is real; the label is wrong. The test needs
      redoing with a power cycle (full report \S7.2).
\end{itemize}

\subsubsection*{The clock-transfer surface}

I tested four frequencies with $\delta$ tuned per cell. The rule I wrote from two
frequencies, ``profile at the highest clock'', is wrong as stated. What holds instead: a
template is bound to its clock inside a band of roughly a factor of two, and the binding
is asymmetric. In all six pairs of distinct clocks, profiling above and attacking below
beats the reverse. Tuning the radius per cell moved the matched 10\,MHz result from
$0.675$ to $0.882$.

\begin{table}[H]
  \centering
  \small
  \begin{tabular}{lrrrr}
\hline
Profiled at & \multicolumn{4}{c}{Attacked at} \\
 & 5\,MHz & 7.5\,MHz & 10\,MHz & 20\,MHz \\
\hline
5\,MHz & \textit{0.884 {\tiny (0.25)}} & 0.706 {\tiny (0.15)} & 0.050 {\tiny (0.60)} & 0.019 {\tiny (0.60)} \\
7.5\,MHz & 0.901 {\tiny (0.25)} & \textit{0.918 {\tiny (0.15)}} & 0.510 {\tiny (0.10)} & 0.038 {\tiny (0.60)} \\
10\,MHz & 0.851 {\tiny (0.10)} & 0.895 {\tiny (0.15)} & \textit{0.882 {\tiny (0.06)}} & 0.132 {\tiny (0.60)} \\
20\,MHz & 0.071 {\tiny (0.60)} & 0.228 {\tiny (0.60)} & 0.305 {\tiny (0.25)} & \textit{0.345 {\tiny (0.60)}} \\
\hline
\end{tabular}

  \caption{Foreground F1 for every combination of profiling and attack clock, with
  $\delta$ tuned for each cell and the chosen value shown.}
  \label{tab:gridbest}
\end{table}

Two further findings on conditions (full report \S\S8.4, 8.7):

\begin{itemize}
\item \textbf{Randomising the accelerator clock is a weak defence.} It costs an attacker
      who profiled at 5\,MHz half his F1, but one who profiled at 10\,MHz keeps $0.737$
      of $0.882$ and still reads three digits in four, because that template sits within
      transfer range of the whole band.
\item \textbf{The template is not what breaks.} In the 5 and 10\,MHz pair, the direction
      that fails retrieves and ranks the true patch \emph{better} than the direction that
      works, at every radius I tested. The loss therefore happens in the greedy assembly
      stage rather than in the leakage or the template match. I cannot yet say what the
      assembly does wrong, and better features would not have helped.
\end{itemize}

\subsection*{4.4 The attacker cannot use arbitrary probe kernels}

I profiled with one-hot probe kernels and attacked traces captured while the FPGA ran the
real trained kernels. The attack collapses: F1 falls to about $0.12$, and pixel accuracy
falls \emph{below} the all-background baseline, $0.811$ against $0.888$, because the
reconstruction paints foreground pixels where only background exists. The recogniser then
reads one image in five. A time gap does not explain this. In \S4.3 a template transferred
across a comparable gap and still reached $0.84$.

\subsection*{4.5 A template does not transfer to a different dataset, and I can say why}

An MNIST-built template attacking Fashion-MNIST scores at chance, $0.034$ above a
label-blind predictor, while a Fashion-built template on the same traces reaches $0.662$.
The template is a dictionary of observed $3\times3$ patches. MNIST covers 156 of the 512
possible patches and Fashion covers 450, with MNIST very nearly a subset, so transfer
works only in the direction of increasing coverage.

Raw F1 does not compare across datasets: the label-blind ceiling moves from $0.204$ on
MNIST to $0.474$ on Fashion. Every cross-dataset number here is reported against its own
chance baseline for that reason.

\begin{table}[H]
  \centering
  \small
  \begin{tabular}{llrrrrr}
\hline
Template & Attacked & $\delta$ & F1 & chance & excess & MSSIM \\
\hline
MNIST & MNIST & 0.060 & 0.882 & 0.204 & 0.852 & 0.852 \\
MNIST & Fashion & 0.060 & 0.432 & 0.412 & 0.034 & 0.432 \\
Fashion & MNIST & 0.040 & 0.810 & 0.204 & 0.762 & 0.761 \\
Fashion & Fashion & 0.020 & 0.801 & 0.412 & 0.662 & 0.626 \\
\hline
\end{tabular}

  \caption{Template dataset against attacked dataset, $\delta$ tuned per cell.
  ``Excess'' is the margin over the label-blind ceiling for that truth set.}
  \label{tab:dsgrid}
\end{table}

\begin{figure}[H]
  \centering
  \includegraphics[width=\textwidth]{figures/fig_dataset_transfer.png}
  \caption{Fashion-MNIST recovered by the template attack.}
  \label{fig:dstransfer}
\end{figure}

\subsection*{4.6 Greyscale: real intensity comes back from power alone}

The template attack classifies 1-bit patches and cannot express a grey level. Removing the
input binarisation and running the generative (Power2Picture-style) route on the
grey-pixel design brings real grey levels back: mean absolute error under $10$ on the
0--255 scale for MNIST, and correlation $0.93$, from a design confirmed bit-exact on
hardware.

One correction. I first credited this to the input width, but a control bitstream shows
that most of the apparent gain came from the wider leakage amplifier in the grey design.
Hold the amplifier equal and grey scores \emph{worse} than binary on foreground F1, better
only on the intensity measures. That narrower claim is the one that survives (full report
\S8.6).

The grey pipeline also works on Fashion-MNIST, a dense non-digit dataset nobody had tested
it on. A Fashion classifier reads the power-recovered images at $0.827$ against $0.830$ on
the originals, so in classification terms the reconstruction loses almost nothing. Pixel
correlation is $0.910$ against $0.929$ in the truth.

\begin{figure}[H]
  \centering
  \includegraphics[width=\textwidth]{figures/fig_fashion_grey.png}
  \caption{Fashion-MNIST greyscale recovered from power alone.}
  \label{fig:fashgrey}
\end{figure}

\subsection*{4.7 The attack does not need the leakage amplifier}

The leakage amplifier is a flop bank I added to the design to make it leak, so a reader
can fairly ask how much of the result belongs to the accelerator and how much to my own
circuit. Read in isolation, the sweep from $4536$ flops down to $72$ looks damaging:
recognition falls $0.857 \to 0.400$ and MSSIM $0.730 \to 0.503$. But every point in that
sweep used the same small profiling budget of about $1500$ traces, which confounds
amplifier size with data volume.

\begin{figure}[H]
  \centering
  \includegraphics[width=\textwidth]{figures/fig_amplifier_sweep.png}
  \caption{The same test images recovered at each amplifier width.}
  \label{fig:ampsweep}
\end{figure}

\begin{table}[H]
  \centering
  \small
  \resizebox{\textwidth}{!}{\begin{tabular}{lrrrrrr}
\hline
Amplifier, profiling set & MAE & Corr. & MSSIM & F1 excess & Recog. & (ceiling) \\
\hline
none ($1$ flop), $56$k traces & 23.98 & 0.878 & 0.724 & 0.632 & 0.746 & 0.898 \\
$72$ flops, $1.5$k traces & 37.26 & 0.750 & 0.503 & 0.422 & 0.400 & 0.830 \\
$72$ flops, $56$k traces & 23.75 & 0.880 & 0.725 & 0.635 & 0.757 & 0.898 \\
$4536$ flops, $1.5$k traces & 20.97 & 0.909 & 0.730 & 0.650 & 0.857 & 0.830 \\
$4536$ flops, $56$k traces & 14.49 & 0.950 & 0.861 & 0.759 & 0.919 & 0.898 \\
\hline
\end{tabular}
}
  \caption{Amplifier size against profiling-set size. Identical loss, batch, epochs,
  split and seed in every cell.}
  \label{tab:ampdata}
\end{table}

Separating the two changes the conclusion. A $63$ times smaller amplifier costs about what
a $37$ times smaller profiling set costs, and either deficit can be partly bought back
with the other. \textbf{The amplifier changes how much profiling data an attacker needs;
it does not decide whether the attack is possible.} A defender should prefer the first
reading, because profiling traces are the cheap resource in this threat model.

Removing the bank altogether leaves one flop registering the output, against $72$ in the
smallest amplified build. That build still reaches MSSIM $0.724$ and recognition $0.746$
on the paper's $60\,000$-trace profiling set, within noise of the $72$-flop build's
$0.725$ and $0.757$. Three earlier claims in this project were wrong about this, and all
three failed the same way, by treating a limitation of my own build as a fact about the
hardware. A combinational path from the convolution datapath to a top-level pad killed the
bank-free builds, which is my mistake and not the board's. One flop fixes it (full report
\S9.7).

\subsection*{4.8 Making the victim more realistic did not weaken the attack}

I cut the window dwell from $8$ cycles to $2$, so the accelerator streams closer to the way
a real one would. The result improved: recognition $0.746 \to 0.807$ and MSSIM
$0.724 \to 0.762$ on the same test images, with the amplifier absent in both cases. That
gap is about six standard errors, so it is not noise, and I had predicted the opposite
before running it. Both bench artifacts I introduced, the leakage amplifier and the slow
dwell, turned out to be things the attack did not need.

\begin{table}[H]
  \centering
  \small
  \resizebox{\textwidth}{!}{\begin{tabular}{lrrrrrr}
\hline
Window dwell & Input dim. & MAE & Corr. & MSSIM & F1 excess & Recog. \\
\hline
$8$ cycles ($32$ samples/window) & 5408 & 23.98 & 0.878 & 0.724 & 0.632 & 0.746 \\
$2$ cycles ($8$ samples/window) & 1352 & 21.44 & 0.902 & 0.762 & 0.681 & 0.807 \\
\hline
\end{tabular}
}
  \caption{Window dwell at zero amplifier, $56$k training traces.}
  \label{tab:dwell}
\end{table}

\subsection*{4.9 Against the published numbers}

Run with Power2Picture's own training protocol, meaning their generator architecture,
optimiser, loss, batch size and $60\,000$-image profiling set, the attack reaches MSSIM
$0.861$, pixel distance $14.5$ and recognition $0.919$ against a ceiling of $0.898$.
Power2Picture report $0.38$, $77$ and $0.53$ on Fashion-MNIST.

\textbf{This is not a fair comparison and I do not present it as one.} They attack a full
quantised LeNet through five on-chip TDC sensors on a shared FPGA, which is a remote attack
against a realistic accelerator. I attack a single convolution layer through an external
shunt. Their measurement is harder in every respect that counts. The comparison
establishes one thing only: this pipeline leaves nothing obvious on the table relative to
the published method.

\section*{5. What I take from these numbers}

\begin{itemize}
\item The attack works on real hardware. In the control condition the digit recogniser
      reads every recovered image correctly.
\item Two causes stacked behind the earlier weak results: the traces came from the wrong
      design, and even correct traces need averaging. Dilution is a physical property of
      this board and will not go away.
\item The fragility has a specific shape. The clock matters, gain does not, dataset
      coverage decides transfer, and the probe kernels must match the deployed ones. The
      original paper never reports what happens when profiling and attack conditions
      differ, which puts this closer to a contribution than to a reproduction.
\item More profiling data does not always help. Template size and search radius have to be
      tuned together, and I reported an untuned radius twice in this project.
\item The open result I find most interesting: when cross-clock transfer fails, the
      template still ranks the true patch well. The greedy assembly stage throws away an
      answer the earlier stages found. That moves the question from the leakage to the
      algorithm.
\end{itemize}

\subsection*{Next steps, in the order I would do them}

\begin{enumerate}
\item Port the on-chip RO/TDC sensor into the grey design. With the amplifier gone, the
      external shunt is the largest remaining deviation from Power2Picture. The RO
      channel's correlation plateaus near $0.33$, so expect the reconstruction to get
      worse rather than better. That is the point, and it is the honest version of the
      threat model.
\item Find out what the greedy assembly does to the cross-condition cells. This is cheap
      and needs no board: dump the candidate lists and replay the assembly. A
      reconstruction that resolves disagreement globally instead of greedily would improve
      on the paper's Algorithm 2 rather than reproduce it.
\item Recapture at 20\,MHz with two samples per cycle instead of four. That separates
      weaker leakage at a high clock from the ChipWhisperer-Lite's front end running at
      80\,MSa/s. Until then I cannot say whether clocking the accelerator faster is a real
      countermeasure, which is the cheapest defence anyone would try.
\item Make the template attack regress a grey patch instead of classifying a binary one,
      which is the missing piece for a direct comparison with the paper.
\item Redo the reprogramming test with a power cycle and put the withdrawn row back with a
      real measurement behind it.
\item Measure transfer across days and temperatures, not only across settings inside one
      session.
\end{enumerate}

\section*{Appendix: reproducing this}

All scripts live in the repository. The analysis scripts do not need the board.

\begin{quote}\footnotesize\ttfamily
host/cw305\_leakage\_capture.py \quad -- capture traces from the CW305\\
attack/hardtests\_20260830/run\_repeat\_sweep.sh \quad -- the averaging sweep\\
attack/hardtests\_20260830/run\_conditions.sh \quad -- the condition transfer tests\\
attack/hardtests\_20260830/run\_probe\_mismatch.sh \quad -- the probe mismatch test\\
attack/hardtests\_20260830/run\_paperscale.sh \quad -- 500-image capture, 300-image template\\
report/catchup\_v2\_2026-08-30/make\_figures.py \quad -- all figures\\
report/catchup\_v2\_2026-08-30/make\_tables.py \quad -- all tables
\end{quote}

Every figure showing a power signal comes from a real capture file, and \texttt{make\_tables.py}
generates the tables from the score files, so nobody typed a number in this report by hand.

\end{document}
